\section{The HEIST Environment \& Benchmark}
\label{sec:env}

We formalize the HEIST benchmark as a Decentralized Partially Observable Markov Decision Process (Dec-POMDP) augmented with asymmetric operational roles and procedural facility generation.

\subsection{Dec-POMDP Formulation}
A mission instance is defined by the tuple:
\begin{equation}
\mathcal{M} = \langle \mathcal{N}, \mathcal{S}, \{\mathcal{A}_i\}, \mathcal{P}, \mathcal{R}, \{\Omega_i\}, \mathcal{O}, \gamma \rangle
\end{equation}
\begin{itemize}[leftmargin=*,nosep]
    \item $\mathcal{N} = \{\text{Scout}, \text{Hacker}, \text{Muscle}, \text{Extractor}\}$ denotes the agent team with $|\mathcal{N}| = 4$.
    \item $\mathcal{S}$ is the global environment state, encompassing the spatial tile grid $G \in \mathbb{Z}^{H \times W}$, agent coordinates $\mathbf{p}_i = (r_i, c_i)$, guard dynamic coordinates and behavioral states, camera scan cones, door states, global alarm $\alpha \in [0, \alpha_{\max}]$, and extraction countdown $\tau_{\text{ext}}$.
    \item $\mathcal{A}_i = \{\text{UP}, \text{DOWN}, \text{LEFT}, \text{RIGHT}, \text{WAIT}, \text{INTERACT}\}$ is the discrete action space ($|\mathcal{A}_i| = 6$). Wall movements are masked out before categorical sampling.
    \item $\Omega_i$ represents the local egocentric observation space. Each agent receives a $7 \times 7$ egocentric grid $o_i$ masked by raytraced fog-of-war, a one-hot role vector $\mathbf{r}_i \in \{0,1\}^4$, an action mask $\mathbf{m}_i \in \{0,1\}^6$, and a 2D mission-orientation compass vector $\mathbf{g}_i \in [-1, 1]^2$ that points toward the agent's current sub-task objective (e.g., the security terminal for the Hacker until terminal override, the vault for the Extractor, and the extraction zone after loot acquisition).
    \item $\gamma = 0.99$ denotes the temporal discount factor.
\end{itemize}

\subsection{Asymmetric Role Specialization}
The four agent roles possess strictly non-interchangeable mechanics (summarized in Table~\ref{tab:stage_specs}):
\begin{itemize}[leftmargin=*,nosep]
    \item \textbf{Scout ($S$):} Possesses an extended sensor perimeter ($r_{\text{vis}} = 8$ vs.\ $r_{\text{vis}} = 3$ for teammates). Using \texttt{INTERACT} within standoff range ($d \le 8$), the Scout tags Points of Interest (POIs: cameras, terminal, vault loot), projecting continuous directional guidance onto teammates' compass vectors $\mathbf{g}_i$.
    \item \textbf{Hacker ($H$):} Bypasses locked security doors via \texttt{INTERACT} and hacks the security terminal through three consecutive interaction turns, permanently deactivating security cameras and unlocking vault access.
    \item \textbf{Muscle ($M$):} Engages and neutralizes patrolling security guards within distance $d \le 2$. Neutralized guards remain incapacitated for the remainder of the episode, preventing alert escalations along escape paths.
    \item \textbf{Extractor ($E$):} Infiltrates the unlocked vault, secures the loot payload ($\$$), and initiates the extraction countdown $\tau_{\text{ext}} = \lfloor 0.5 \times T_{\max} \rfloor$, requiring all four agents to reach the extraction zone.
\end{itemize}

\subsection{Security Dynamics and Alarm Mathematics}
Security guards operate under a three-state finite-state machine: \emph{Patrol} (waypoint navigation), \emph{Investigate} (pathfinding to detected agent locations), and \emph{Search} (breadth-first search within radius $5$). If any guard closes to $\text{CATCH\_DISTANCE} = 1$, the episode terminates in immediate defeat.

Global alarm $\alpha$ accumulates continuously based on security violations:
\begin{equation}
\Delta \alpha = \mathbb{I}(\text{terminal active}) \cdot \sum_{c \in \text{Cameras}} \alpha_{\text{cam}} \cdot N_{\text{visible}}(c) + \alpha_{\text{events}}
\end{equation}
where event increments include guard line-of-sight exposure ($+10.0$), continuous tracking ($+1.5$/step), guard neutralizations ($+5.0$), door bypasses ($+3.0$), and extraction timeout ($+15.0$). If $\alpha \ge \alpha_{\max}$, facility lockdown is triggered, resulting in immediate failure ($R_{\text{team}} = -10.0$).

\subsection{Phased Potential-Based Reward Shaping}
The team objective reward provides $+15.0$ upon successful 4-agent extraction with loot and $-10.0$ on catastrophic failure, alongside step bleed $r_{\text{bleed}} = -2.0 / T_{\max}$. To mitigate credit assignment starvation across long planning horizons without altering the optimal policy~\cite{ng1999policy}, we introduce role-phased potential-based reward shaping:
\begin{equation}
\Phi_i(s) = -\|\mathbf{p}_i - \mathbf{p}_{\text{target}}\|_1, \quad r_{i,\text{shaping}} = \gamma \Phi_i(s') - \Phi_i(s)
\end{equation}
Potential functions activate sequentially per role upon respective subtask completion: Hacker after terminal override, Muscle after guard clearance, Scout upon POI tagging, and Extractor upon loot acquisition.
